\begin{figure}[t]
    \centering
    \resizebox{\textwidth}{!}{
        \Timeline[1905,2020,0.13]{
            \event{1911}{thorndike1911law}{\textbf{Law of Effect}, learning driven by the consequences of actions.}{5cm}{0.2cm}{0.05cm}{\scriptsize}{blue!10}
            \event{1938}{skinner1938behavior}{\textbf{Operant Conditioning}, behavior shaped by rewards and punishments.}{5cm}{1cm}{-0.04cm}{\scriptsize}{cyan!10}
            \event{1948}{turing1948intelligent}{\textbf{Learning Machines}, machines learning through reward-like feedback.}{5cm}{1cm}{0.12cm}{\scriptsize}{gray!15}
            \event{1958}{bellman1957dynamic,howard1960mdp}{\textbf{Dynamic Programming and Markov Decision Processes}, the first mathematical framework for sequential decision-making.}{5.5cm}{1cm}{0.05cm}{\scriptsize}{red!10}
            \event{1985}{BartoSuttonAnderson1983,sutton1988td}{\textbf{Actor-Critic and Temporal-Difference Learning}, incremental learning directly from experience.}{5.5cm}{1cm}{-0.04cm}{\scriptsize}{green!10}
            \event{1989}{watkins1989qlearning,watkins1992qlearning}{\textbf{Q-Learning}, optimal action-values learned without a model of the environment.}{5cm}{1cm}{-0.1cm}{\scriptsize}{yellow!15}
            \event{1997}{tsitsiklis1996analysis,suttonbarto1998}{Convergence analysis of \textbf{function approximation} and publication of the reference book \textbf{Reinforcement Learning: An Introduction}.}{5.5cm}{1cm}{0.12cm}{\scriptsize}{purple!10}
            \event{2013}{mnih2013playingatarideepreinforcement}{\textbf{Deep Q-Networks}, deep networks approximate value functions at scale.}{5cm}{1cm}{0.05cm}{\scriptsize}{orange!10}
        }
    }
    \caption{Evolution of reinforcement learning prior to its fusion with deep learning.}
    \label{fig:evolution_reinforcement_learning_classical}
\end{figure}
